%%%%%%%%%%  scan idx 315  |  folio 302  %%%%%%%%%%
\SGAPrintPage{302}%
\(\theta=\Phi(1+Xt)\);
let us show that \(\theta\) satisfies conditions i) and ii).

i): \(H(f_1)(\theta)=H(f_1)\circ\Phi(1+Xt)=\Phi\circ\Ghat(f_1)(1+Xt)=\Phi(1)=0\).

ii): \(u_H^n(\theta)=\Phi(1+\ldots+X_nt^n)\) by (1.6.1); now \(1+\ldots+X_nt^n\equiv 1+\ldots+X_{n-1}t^{n-1}\ \mathrm{mod.}\ F^n\Ghat\), and \(1+\ldots+X_{n-1}t^{n-1}=\Ghat(f_n)(1+\ldots+X_nt^n)\). Since \(\Phi\circ\Ghat(f_n)=H(f_n)\circ\Phi\), and since \(\Phi\) is a morphism of filtered functors, we obtain ii).

Finally, we must show that if \(\theta\in H(K[X])\) satisfies i) and ii), there exists \(\Phi\in\Hom(\Ghat,H)\) such that \(\theta=\Phi(1+Xt)\).

We first determine the value of \(\Phi\) on the element \(1+\ldots+X_nt^n\) of \(\Ghat(K[X_1,\ldots,X_n])\) by means of (1.6.1); thus, by definition, \(\Phi(1+\ldots+X_nt^n)=\sum\limits_{i=1}^n \Phi(1+T_it)\), the \(\Phi(1+T_it)\) being known by functoriality through the morphisms \(u_i^n\).

We then glue the \(\Phi(1+\ldots+X_nt^n)\) for varying \(n\) by means of the commutative diagram:
\sgaStableTarget{sga6.e05.eq.1-6-2}\[
\begin{tikzcd}
K[X_1,\ldots,X_n] \arrow[r,"f_n"] \arrow[d,"\sigma_n"'] &
K[X_1,\ldots,X_n] \arrow[d,"\sigma'_n"']\\
K[T_1,\ldots,T_n] \arrow[r,"f'_n"] &
K[T_1,\ldots,T_n],
\end{tikzcd}
\tag{1.6.2}
\]
where \(f'_n\) is determined by \(f'_n(T_i)=T_i\) if \(i<n\), \(f'_n(T_n)=0\), where \(\sigma_n\) sends \(X_i\) to the \(i\)-th symmetric function of the \(T_j\), and \(\sigma'_n\) sends \(X_i\) to the \(i\)-th symmetric function of \(T_1,\ldots,T_{n-1}\) for \(i<n\), and \(X_n\) to \(0\).

Indeed, this gives:
\[
\begin{aligned}
H(f_n)\bigl(\Phi(1+\ldots+X_nt^n)\bigr)
&=H(f'_n)\left(\sum_{i=1}^n \Phi(1+T_it)\right)
=\sum_{i=1}^{n-1}\Phi(1+T_it)\ \text{(by 1.6.i))}\\
&=\Phi(1+\ldots+X_{n-1}t^{n-1})
\end{aligned}
\]
%%%%%%%%%%  scan idx 316  |  folio 303  %%%%%%%%%%
\SGAPrintPage{303}%
(working in \(H(K[T_1,\ldots,T_{n-1}])\) by means of the identifications 1.4).

In view of 1.6.ii), \(\Phi(1+\ldots+X_nt^n)\equiv\Phi(1+\ldots+X_{n-1}t^{n-1})\ \mathrm{mod.}\ F^nH(K[X_i]_{i\geq1})\), and the \(\Phi(1+\ldots+X_nt^n)\) define an element \(\psi\) of \(H(K[X_i]_{i\geq1})\), since this is a separated complete group; we therefore put \(\Phi(\varphi)=\psi\).

Finally, let \(A\in\operatorname{Ob}(\mathcal C)\), and \(\chi\in\Ghat(A)\); we put \(\Phi(\chi)=H(f_\chi)(\psi)\), where \(f_\chi\) is the homomorphism from \(K[X_i]_{i\geq 1}\) to \(A\) defined in 1.2. Since we have already defined \(\Phi\) on the elements \(1+\ldots+X_nt^n=\varphi_n\), we must verify that the two definitions give the same result, that is:
\[
  H(f_{\varphi_n})(\psi)=\Phi(1+\ldots+X_nt^n).
\]
Now, for every \(p\), we have: \(\psi\equiv \Phi(1+\ldots+X_pt^p)\ \mathrm{mod.}\ F^{p+1}H(K[X_i]_{i\geq 1})\). Hence, for every \(p\geq n\); \(H(f_{\varphi_n})(\psi)\equiv H(f_{\varphi_n})\bigl(\Phi(1+\ldots+X_pt^p)\bigr)\ \mathrm{mod.}\ F^{p+1}H(K[X_i]_{i\geq 1})\). Using the identifications of 1.4 and the gluing of the \(\Phi(\varphi_m)\) for varying \(m\), we obtain: \(H(f_{\varphi_n})(\Phi(\varphi_p))=\Phi(\varphi_n)\).
Thus, finally:
\[
  \forall p\geq n,\quad
  H(f_{\varphi_n})(\psi)\equiv \Phi(1+\ldots+X_nt^n)
  \ \mathrm{mod.}\ F^pH(K[X_i]_{i\geq 1}).
\]
Since \(H\) is separated, we obtain the required equality.

We have thus defined a system of maps \(\Phi_A:\Ghat(A)\to H(A)\); to verify that they constitute a morphism of functors (with values in \(\underline{\mathrm{Ens}}\)), we must show that for every morphism \(h:A\to A'\), the diagram:
\[
\begin{tikzcd}
\Ghat(A) \arrow[r,"\Ghat(h)"] \arrow[d,"\Phi_A"'] &
\Ghat(A') \arrow[d,"\Phi_{A'}"]\\
H(A) \arrow[r,"H(h)"] & H(A')
\end{tikzcd}
\]
%%%%%%%%%%  scan idx 317  |  folio 304  %%%%%%%%%%
\SGAPrintPage{304}%
is commutative. Let \(\chi\in\Ghat(A)\); then: \(h\circ f_\chi=f_{\Ghat(h)(\chi)}\); hence:
\[
  H(h)(\Phi_A(\chi))=H(h)\circ H(f_\chi)(\psi)
  =H(f_{\Ghat(h)(\chi)})(\psi)=\Phi_{A'}(\Ghat(h)(\chi)).
\]

It therefore remains to show that the morphism of functors thus constructed is compatible with the structures of filtered groups.

Suppose that \(\chi\in F^n\Ghat(A)\). To say that \(\chi\) has filtration \(\geq n\) is equivalent to saying that \(f_\chi:K[X_i]_{i\geq 1}\to A\) factors through the endomorphism \(g_n\) of \(K[X_i]_{i\geq 1}\) defined by: \(g_n(X_i)=0\) if \(1\leq i\leq n-1\), and \(g_n(X_i)=X_i\) if \(i\geq n\).
To show that \(\Phi\) is compatible with the filtrations, it therefore suffices to show that \(H(g_n)(\Phi(\varphi))\in F^nH\). Now, by construction:
\[
  \Phi(\varphi)\equiv \Phi(1+\ldots+X_{n-1}t^{n-1})\ \mathrm{mod.}\ F^nH,
\]
hence:
\[
  H(g_n)(\Phi(\varphi))
  \equiv H(g_n)\bigl(\Phi(1+\ldots+X_{n-1}t^{n-1})\bigr)
  =0\ \mathrm{mod.}\ F^nH.
\]

To verify the additivity of \(\Phi\), it suffices to do this for the elements of finite degree of \(\Ghat(A)\), and then pass to the limit. Thus let \(\chi_1,\chi_2\) be two elements of degree \(n\) of \(\Ghat(A)\), and \(\chi\) their product. Consider the commutative diagram:
\[
\begin{tikzcd}[column sep=tiny]
K[X_1,\ldots,X_{2n}] \arrow[rr,"f_\chi"] \arrow[d,hook] && A\\
\bigl(K[T_1,\ldots,T_{2n}]\bigr)^{\mathfrak S_n\times\mathfrak S_n}
\arrow[urr,"f'_\chi"] && ,
\end{tikzcd}
\]
where the arrows are defined as follows: \(\mathfrak S_n\times\mathfrak S_n\) acts on \(K[T_1,\ldots,T_{2n}]\)
%%%%%%%%%%  scan idx 318  |  folio 305  %%%%%%%%%%
\SGAPrintPage{305}%
in such a way that the first factor permutes \(T_1,\ldots,T_n\) and the second \(T_{n+1},\ldots,T_{2n}\); \(K[X_1,\ldots,X_{2n}]\) is embedded in the ring of invariants by sending \(X_j\) to the \(j\)-th symmetric function of the \(T_i\); \(f'\) sends the \(i\)-th symmetric function of \(T_1,\ldots,T_n\) (resp. \(T_{n+1},\ldots,T_{2n}\)) to the \(i\)-th coefficient of \(\chi_1\) (resp. \(\chi_2\)). Then:
\[
\begin{aligned}
\Phi(\chi)
&=H(f_\chi)\left(\Phi\left(\prod_{i=1}^{2n}(1+T_it)\right)\right)
=H(f'_\chi)\left(\Phi\left(\prod_{i=1}^{2n}(1+T_it)\right)\right);\\
\Phi(\chi_1) &=H(f'_\chi)\left(\Phi\left(\prod_{i=1}^{n}(1+T_it)\right)\right)\quad(\text{resp. }\chi_2,\ i=n+1,\ldots,2n).
\end{aligned}
\]
Since \(H(f'_\chi)\) is a homomorphism of groups, we are reduced to the case where \(\chi=\prod_{i=1}^{2n}(1+T_it)\), and in this case the additivity of \(\Phi\) follows from the definition of the \(\Phi(1+X_1t+\ldots+X_k t^k)\).

\sgaStableTarget{sga6.e05.para.1-7}1.7. We can generalize the preceding considerations to a product of functors \(\Ghat\); the functor \((\Ghat)^k\) is representable by the ring
\(K[(X_i^{(j)})]\), \(j=1,\ldots,k\), and \(i\geq 1\), where \(X_i^{(j)}\) has weight \(i\). We denote by \(\varphi^{(j)}\) the series \(1+X_1^{(j)}t+\ldots+X_n^{(j)}t^n+\ldots\); the element \((\varphi^{(j)})_j\) of \(\bigl(\Ghat(K[(X_i^{(j)})])\bigr)^k\) has the universal property analogous to that of \(\varphi\) (cf. 1.2).

We can introduce the morphisms \(f_n^{(j)}\), \(u_H^{n,(j)}\) as before, by repeating the construction for each family \((X_i^{(j)})_{i\geq 1}\) separately. We obtain a result analogous to Theorem 1.6 by seeking the homomorphisms of functors from \((\Ghat)^k\) to \(H\) that are \(\mathbb Z\)-multilinear and compatible with the filtrations in each of the \(k\) arguments separately: they are determined by their value on the element
%%%%%%%%%%  scan idx 319  |  folio 306  %%%%%%%%%%
\SGAPrintPage{306}%
\((1+X^{(j)}t)_{j=1,\ldots,k}\) of \(\bigl(\Ghat(K[(X^{(j)})_j])\bigr)^k\); an element \(\theta\) of \(H(K[(X^{(j)})_j])\) determines a morphism of functors if and only if it satisfies:
\[
\begin{array}{ll}
\text{i)}& \forall j,\quad H(f_1^{(j)})(\theta)=0\\[0.3em]
\text{ii)}& \forall n,j,\quad
H(f_n^{(j)})\bigl(u_H^{n,(j)}(\theta)\bigr)
\equiv u_H^{n,(j)}(\theta)\ \mathrm{mod.}\ F^nH.
\end{array}
\]

The reader will readily verify that, apart from complications of notation, it suffices to copy the proof of 1.6.

\sgaStableTarget{sga6.e05.para.1-8}1.8. We denote by \(\mathcal C_0\) the category of \(K\)-algebras, and by
\[
  i:\mathcal C_0\longrightarrow \mathcal C,
\]
the functor that assigns to a \(K\)-algebra \(A\) the \(K\)-algebra \(A\) endowed with the coarse filtration; we denote by
\[
  g:\mathcal C\longrightarrow\mathcal C_0,
\]
the functor ``forget the filtration.'' These functors satisfy the following relations:
\[
  g\circ i=\operatorname{Id}_{\mathcal C_0};\qquad
  \Ghat\circ i\circ g=\Gg;\qquad
  \Gg\circ i\circ g=\Gg
\]
(cf. 1.3, ex. 2)

\sgaStableTarget{sga6.e05.lem.1-9}\textbf{Lemma 1.9.} \uline{Let} \(H:\mathcal C_0\to\underline{\mathrm{Ab}}\,\underline{\mathrm{fil}}\) \uline{be a functor on} \(\mathcal C_0\); \uline{then the canonical homomorphism}:
\[
  \Hom_{\underline{\mathrm{Ab}}\,\underline{\mathrm{fil}}}(\Ghat,H\circ g)
  \longrightarrow
  \Hom_{\underline{\mathrm{Ab}}\,\underline{\mathrm{fil}}}(\Ghat\circ i,H)
\]
